\chapter{Multiplicity, Entropy, and Constrained Equilibrium}

\lectureoverview{A probability distribution becomes a thermodynamic object
once its states carry energies. We first connect entropy increase to
equilibration and use \(dE=T\,dS\) to derive the direction of heat flow. We
then replace one system by a large collection of replicas, count the
arrangements associated with each set of occupation numbers, and show that
the logarithm of this multiplicity is \(N\) times the entropy. Equilibrium is
therefore a constrained maximum-entropy problem. Lagrange multipliers provide
the method needed to solve it.}

\section{The Three Basic Quantities}

Let \(i\) label the discrete states available to a system. For the moment,
discreteness is useful because it lets us count states without introducing
phase-space measures. If the probability of state \(i\) is \(P_i\), then
\begin{equation}
 \sum_i P_i=1.
 \label{eq:sm3-normalization}
\end{equation}
If the energy of that state is \(E_i\), the mean energy is
\begin{equation}
 \overline E=\sum_iP_iE_i.
 \label{eq:sm3-mean-energy}
\end{equation}
The word \emph{mean} matters. The system need not possess the value
\(\overline E\) in any one microscopic state; it is the probability-weighted
average over the possible values.

\begin{classroomqa}
\classquestion{Is there a useful distinction here between probability and
statistics?}
\lecturerresponse{Probability supplies a model for possible outcomes, while
statistics usually begins with observed data and asks what can be inferred
from it. The present calculation uses probability distributions, although the
same averages naturally appear in statistical data. The informal exchange in
class treats the words lightly because that distinction does not affect the
physics that follows.
\lecturetimestamp{00:01:39--00:02:21}}
\end{classroomqa}

The entropy of the distribution is
\begin{equation}
 S[P]=-\sum_iP_i\log P_i.
 \label{eq:sm3-entropy}
\end{equation}
This expression is nonnegative. Every nonzero probability obeys
\(0<P_i\leq1\), so \(\log P_i\leq0\), and the conventional limiting value
\(P_i\log P_i\to0\) applies when \(P_i\to0\).

\begin{figure}[t]
\centering
\includegraphics[width=0.88\textwidth]{lecture_03_figure_02.png}
\caption{The opening board collects the entropy, normalization, and
mean-energy formulas together with a qualitative family of equilibrium
distributions.}
\lectureframe{00:06:44}
\end{figure}

\section{Equilibrium as a Family of Distributions}

The previous lecture left us with a one-parameter family of equilibrium
distributions. One may label that family by temperature, but for now it is
more useful to label it by the mean energy \(\overline E\). At the lowest
possible energy the distribution is concentrated on a ground state:
\[
 P_i=\delta_{i0}.
\]
There is then no uncertainty about the state, and \(S=0\). As the equilibrium
distribution spreads over more energy levels, both its mean energy and its
entropy rise.

This is a statement about the thermal family, not about every distribution
that can be invented. An arbitrary family could shift toward larger energies
while becoming narrower. The equilibrium distributions to be derived do not
behave that way in the ordinary positive-temperature regime. Along that
family,
\begin{equation}
 \frac{dS}{d\overline E}>0.
 \label{eq:sm3-positive-temperature-branch}
\end{equation}
Systems with bounded spectra can also have branches on which this derivative
is negative. Those negative-temperature regimes are real, but they are set
aside during the present argument.

Thermal equilibrium means more than microscopic stillness. Molecules continue
to move, collide, and exchange energy. What has become stationary is the
macroscopic description: the distribution no longer exhibits a net drift,
and separated portions of the system no longer sustain a net energy flow.

\section{The Thermodynamic Laws in Statistical Language}

The first law is energy conservation. The second law says that the
thermodynamic entropy of an isolated macroscopic system does not decrease:
\begin{equation}
 \Delta S_{\mathrm{tot}}\geq0.
 \label{eq:sm3-second-law}
\end{equation}
An equality describes equilibrium or an ideal reversible change; a spontaneous
relaxation ordinarily gives a strict increase.

In the probabilistic picture, relaxation tends to spread weight over the
macroscopic alternatives that we still distinguish. This statement requires
the qualification emphasized in the lecture: we are \emph{coarse-graining}.
Exact Hamiltonian evolution preserves fine-grained phase-space entropy.
Macroscopic entropy grows because we cease to track microscopic correlations
and details. Losing such information broadens the reduced description. A
spontaneous narrowing would instead mean that we had learned microscopic
information without performing a measurement.

The zeroth law asserts that systems can reach thermal equilibrium and that
equilibrium is transitive. If \(A\) is in equilibrium with \(B\), and \(B\)
with \(C\), then \(A\) is in equilibrium with \(C\). Temperature supplies the
number that makes this transitivity transparent:
\[
 T_A=T_B,\qquad T_B=T_C
 \quad\Longrightarrow\quad
 T_A=T_C.
\]
Out of equilibrium, energy flows from the subsystem with larger temperature
to the subsystem with smaller temperature. At equal temperatures the net flow
ceases. We now derive that operational statement from the entropy definition
of temperature.

\section{Why Energy Flows from Hot to Cold}

Consider two weakly coupled subsystems \(A\) and \(B\). Each is internally
equilibrated, but their temperatures differ. Choose the labels so that
\begin{equation}
 T_B>T_A>0.
 \label{eq:sm3-temperature-order}
\end{equation}
During a small exchange, conservation of the combined energy gives
\begin{equation}
 dE_A+dE_B=0.
 \label{eq:sm3-energy-conservation}
\end{equation}
The second law gives
\begin{equation}
 dS_A+dS_B\geq0.
 \label{eq:sm3-entropy-change}
\end{equation}
For each internally equilibrated subsystem, temperature is defined by
\begin{equation}
 dE_X=T_X\,dS_X,
 \qquad X=A,B.
 \label{eq:sm3-temperature-differential}
\end{equation}
Substitution into \eqref{eq:sm3-energy-conservation} yields
\begin{equation}
 T_A\,dS_A+T_B\,dS_B=0,
\end{equation}
and hence
\begin{equation}
 dS_B=-\frac{T_A}{T_B}\,dS_A.
 \label{eq:sm3-dsb}
\end{equation}
Putting this result into \eqref{eq:sm3-entropy-change},
\begin{align}
 dS_{\mathrm{tot}}
 &=dS_A-\frac{T_A}{T_B}dS_A\\
 &=\frac{T_B-T_A}{T_B}\,dS_A\geq0.
 \label{eq:sm3-flow-entropy}
\end{align}
Because \(T_B-T_A>0\) and \(T_B>0\), a spontaneous exchange requires
\[
 dS_A>0.
\]
Equation \eqref{eq:sm3-temperature-differential} then gives
\begin{equation}
 dE_A=T_A\,dS_A>0,
 \qquad
 dE_B=-dE_A<0.
 \label{eq:sm3-hot-to-cold}
\end{equation}
Energy has flowed from \(B\), the hotter subsystem, to \(A\), the colder one.
Here heat is simply energy in transit. The flow continues until
\(T_A=T_B\), when the first-order entropy change no longer favors either
direction.

The same result can be displayed in one line:
\begin{equation}
 dS_{\mathrm{tot}}
 =dE_A\left(\frac{1}{T_A}-\frac{1}{T_B}\right)\geq0.
 \label{eq:sm3-compact-flow}
\end{equation}
The sign of the inverse-temperature difference fixes the sign of \(dE_A\).
This compact form will later fit naturally with the canonical ensemble.

\begin{classroomqa}
\audiencequestion{Does this argument assume that \(A\) and \(B\) are each
already in internal equilibrium?}
\lecturerresponse{Yes. Assigning a single \(T_A\) or \(T_B\) makes that
assumption. If a subsystem is not internally equilibrated, divide it into
smaller regions and apply the argument locally wherever a temperature can be
defined. The global relaxation is then assembled from exchanges among those
regions.
\lecturetimestamp{00:24:38--00:25:22}}
\end{classroomqa}

\section{What Counts as a State?}

Before asking for the equilibrium probabilities, we must know what the label
\(i\) means. A microscopic state is the most complete specification of the
system allowed by the underlying laws of physics. In classical mechanics,
one pictures the positions and momenta of all particles. In quantum
mechanics, a state is instead a complete quantum state; the uncertainty
principle prevents a simultaneous sharp specification of every classical
position and momentum.

For a quantum system in a finite box, the relevant state counting can be
discrete. Classical statistical mechanics emerges in an appropriate limit,
where sums over densely spaced quantum states become phase-space integrals.
That limiting process also introduces the phase-space cell needed to make
continuous entropy dimensionless; it is not obtained merely by replacing a
summation sign with an integral sign.

\begin{classroomqa}
\audiencequestion{For a gas, should a state be imagined as the position and
momentum of every molecule?}
\lecturerresponse{That is the correct classical picture. More generally, a
state means everything that could in principle be known about the system,
subject to the laws of physics. Quantum mechanics changes the form of that
maximal description and makes the finite-volume counting discrete. Recovering
the classical integral from the quantum sum is a genuine limiting argument,
not a one-line substitution.
\lecturetimestamp{00:25:31--00:27:38}}
\end{classroomqa}

\section{A Heat Bath Made of Replicas}

The hard question is now unavoidable: after equilibrium has been established,
what determines \(P_i\)?

Place the subsystem in contact with a very large heat bath. Energy may pass
back and forth, but the bath has so many degrees of freedom that the transfer
of a small amount of energy barely changes its temperature. A particularly
convenient bath consists of many weakly coupled replicas of the subsystem
itself. One replica is the system under study; the others constitute its
environment.

This construction often mirrors a real macroscopic system. Divide a large
nearly homogeneous body into many similar cells. Each cell is a subsystem,
and the remaining cells act as its bath. Denote the total number of replicas
by \(N\). The classroom demonstration uses nine replicas and three colored
state labels, but the counting is independent of those particular numbers.

Let \(n_i\) be the number of replicas in state \(i\). These are the
\emph{occupation numbers}. There may be infinitely many possible energy
levels, but at finite total energy most very high levels have negligible or
zero occupation.

The first constraint is simply the number of replicas:
\begin{equation}
 \sum_i n_i=N.
 \label{eq:sm3-number-constraint}
\end{equation}
The second is the total energy:
\begin{equation}
 \sum_i n_iE_i=N\overline E.
 \label{eq:sm3-total-energy-constraint}
\end{equation}
The right-hand side is extensive. If the number of replicas is doubled while
their statistical state is held fixed, the total energy doubles, whereas the
mean energy per replica remains \(\overline E\).

\begin{classroomqa}
\audiencequestion{Must the replica construction also satisfy an energy
constraint?}
\lecturerresponse{Yes. The occupation numbers must add to \(N\), and their
energy-weighted sum must equal the fixed total energy \(N\overline E\).
Those are precisely the two constraints that will later restrict the entropy
maximum.
\lecturetimestamp{00:37:12--00:39:54}}
\end{classroomqa}

The probability that a randomly selected replica occupies state \(i\) is its
relative frequency:
\begin{equation}
 P_i=\frac{n_i}{N},
 \qquad
 n_i=NP_i.
 \label{eq:sm3-frequency-probability}
\end{equation}

\begin{figure}[t]
\centering
\includegraphics[width=0.82\textwidth]{lecture_03_figure_03.png}
\caption{The energy-level picture and the identification of probability with
the occupation fraction \(P_i=n_i/N\).}
\lectureframe{00:41:40}
\end{figure}

\begin{classroomqa}
\audiencequestion{Is the notation \(N(i)\) on the board meant to differ from
\(n_i\)?}
\lecturerresponse{No. Both denote the occupation number of state \(i\). We
use \(n_i\) consistently and reserve capital \(N\) for the total number of
replicas.
\lecturetimestamp{00:41:28--00:41:33}}
\end{classroomqa}

Dividing \eqref{eq:sm3-number-constraint} and
\eqref{eq:sm3-total-energy-constraint} by \(N\) recovers the probability
constraints:
\begin{equation}
 \sum_iP_i=1,
 \qquad
 \sum_iP_iE_i=\overline E.
 \label{eq:sm3-probability-constraints}
\end{equation}
The ensemble frequency interpretation has therefore reproduced the equations
with which we began.

\section{Multiplicity of an Occupation Pattern}

Fix a set of occupation numbers \(\{n_i\}\). How many assignments of the
\(N\) distinguishable replicas to the states realize exactly that set? The
answer is the multinomial coefficient
\begin{equation}
 C(\{n_i\})=\frac{N!}{\prod_i n_i!}.
 \label{eq:sm3-multiplicity}
\end{equation}
To see the denominator, first arrange all \(N\) labeled replicas in \(N!\)
orders. Permuting the \(n_i\) replicas that occupy the same state does not
produce a new occupation assignment, so divide by \(n_i!\) for every state.

\begin{figure}[t]
\centering
\includegraphics[width=0.88\textwidth]{lecture_03_figure_04.png}
\caption{The multinomial multiplicity
\(C=N!/\prod_i n_i!\) written on the board.}
\lectureframe{00:46:55}
\end{figure}

Unoccupied levels cause no problem because \(0!=1\). If all replicas occupy
one state, then
\[
 C=\frac{N!}{N!}=1,
\]
as expected. If three replicas occupy three distinct specified states, then
\[
 C=\frac{3!}{1!\,1!\,1!}=6.
\]
With two replicas occupying two specified states, the formula gives two
assignments: exchange which labeled replica occupies which state.

\begin{classroomqa}
\audiencequestion{Can more than one set of occupation numbers satisfy the two
constraints?}
\lecturerresponse{Usually many sets do. Equation
\eqref{eq:sm3-multiplicity} first counts the arrangements for one fixed set
\(\{n_i\}\), without yet deciding whether it satisfies the constraints. The
next stage compares all constraint-compatible sets and finds which one has the
largest multiplicity.
\lecturetimestamp{00:48:07--00:48:43}}
\end{classroomqa}

The statistical assumption is that the microscopic assignments compatible
with the isolated setup are treated symmetrically. A macroscopic occupation
pattern realized by more assignments is therefore more probable. For large
\(N\), the distribution over the fractions \(n_i/N\) becomes sharply peaked.
The equilibrium fractions are the values that maximize \(C\), subject to
\eqref{eq:sm3-number-constraint} and
\eqref{eq:sm3-total-energy-constraint}.

\section{Stirling's Approximation}

Factorials become enormous long before \(N\) reaches a macroscopic value, so
we work with their logarithms. Begin with
\begin{equation}
 \log N!=\sum_{x=1}^{N}\log x.
 \label{eq:sm3-log-factorial-sum}
\end{equation}
For large \(N\), the sum is approximated by the area under \(\log x\):
\begin{align}
 \log N!
 &\simeq\int_1^N\log x\,dx\\
 &=\left[x\log x-x\right]_1^N\\
 &=N\log N-N+1.
 \label{eq:sm3-stirling-integral}
\end{align}
The antiderivative follows immediately from
\[
 \frac{d}{dx}(x\log x-x)=\log x.
\]
At leading extensive order, the additive constant may be dropped:
\begin{equation}
 \log N!\simeq N\log N-N.
 \label{eq:sm3-stirling-log}
\end{equation}
Exponentiating gives the simplified form used in the lecture,
\begin{equation}
 N!\simeq N^N e^{-N}.
 \label{eq:sm3-stirling-simple}
\end{equation}

\begin{figure}[t]
\centering
\includegraphics[width=0.88\textwidth]{lecture_03_figure_05.png}
\caption{The leading Stirling approximation and its derivation by replacing
the sum of \(\log x\) with an integral.}
\lectureframe{01:03:45}
\end{figure}

\begin{classroomqa}
\audiencequestion{Does evaluating \(x\log x-x\) at the lower limit leave an
extra \(+1\)?}
\lecturerresponse{Yes. Equation \eqref{eq:sm3-stirling-integral} contains that
constant. It is negligible compared with the terms proportional to \(N\) as
\(N\to\infty\), which is why the leading formula omits it. The sum itself
must begin at \(x=1\); \(\log0\) is not defined.
\lecturetimestamp{01:04:12--01:04:35}}
\end{classroomqa}

\begin{classroomqa}
\audiencequestion{Is a factor proportional to \(\sqrt{N}\) also missing?}
\lecturerresponse{Yes. The sharper asymptotic formula is
\[
 N!\sim\sqrt{2\pi N}\left(\frac{N}{e}\right)^N.
\]
Its logarithm adds only \(\tfrac12\log(2\pi N)\), which is subextensive
compared with \(N\log N-N\). It matters for accurate finite-\(N\)
probabilities but not for the leading entropy calculation of a macroscopic
system.
\lecturetimestamp{01:04:35--01:05:40}}
\end{classroomqa}

Stirling's approximation must also be used with care in the denominator of
\eqref{eq:sm3-multiplicity}. In the thermodynamic limit the states carrying
nonzero limiting probabilities have \(n_i=NP_i\), so their occupations grow
with \(N\). States with tiny or vanishing occupation can be handled by a
finite-energy cutoff and then restored by a limit. The leading extensive
entropy is unaffected by the omitted subleading terms.

\section{The Logarithm of Multiplicity Is Entropy}

Apply \eqref{eq:sm3-stirling-simple} to the numerator and all macroscopically
occupied denominator factors:
\begin{equation}
 C\simeq
 \frac{N^Ne^{-N}}
 {\prod_i n_i^{\,n_i}e^{-n_i}}.
 \label{eq:sm3-c-stirling}
\end{equation}
Since \(\sum_i n_i=N\), the leading exponential factors cancel:
\begin{equation}
 C\simeq\frac{N^N}{\prod_i n_i^{\,n_i}}.
 \label{eq:sm3-c-leading}
\end{equation}
It is easier to maximize a logarithm than a product, and the logarithm is
monotone, so both maxima occur at the same occupation pattern:
\begin{equation}
 \log C=N\log N-\sum_i n_i\log n_i.
 \label{eq:sm3-logc-occupations}
\end{equation}
Now substitute \(n_i=NP_i\):
\begin{align}
 \log C
 &=N\log N-\sum_iNP_i\log(NP_i)\\
 &=N\log N-N\sum_iP_i\log N
   -N\sum_iP_i\log P_i\\
 &=-N\sum_iP_i\log P_i,
 \label{eq:sm3-logc-entropy-derivation}
\end{align}
where normalization cancels the two \(N\log N\) terms. Thus, to leading
order in \(N\),
\begin{equation}
 \boxed{\log C=NS}.
 \label{eq:sm3-logc-ns}
\end{equation}

\begin{figure}[t]
\centering
\includegraphics[width=0.88\textwidth]{lecture_03_figure_06.png}
\caption{The board calculation connecting
\(\log C=N\log N-\sum_i n_i\log n_i\) to
\(-N\sum_iP_i\log P_i=NS\).}
\lectureframe{01:16:10}
\end{figure}

This relation is the conceptual center of the chapter. Entropy is not an
unrelated quantity imposed on the counting problem. It is the multiplicity
per replica, measured logarithmically. Maximizing the number of microscopic
arrangements is therefore identical, at leading order, to maximizing the
entropy of one replica.

\begin{classroomqa}
\audiencequestion{Would the multiplicity not be largest when all occupation
numbers are equal?}
\lecturerresponse{It is uniform when no state is distinguished and only
\(\sum_i n_i=N\) is imposed over a finite set of states. The fixed-energy
constraint changes the answer because the values \(E_i\) are different.
Equal occupations generally fail to produce the specified mean energy.
\lecturetimestamp{01:17:07--01:18:20}}
\end{classroomqa}

\begin{classroomqa}
\audiencequestion{Capital \(N\) is large, but why may Stirling's approximation
be applied to every \(n_i\) in the denominator?}
\lecturerresponse{For each state with a nonzero limiting probability,
\(n_i=NP_i\) grows proportionally to \(N\). One may first retain a finite
energy range, take the large-\(N\) limit there, and then enlarge the range.
Sparse tail states require separate finite corrections, but they do not alter
the leading extensive entropy.
\lecturetimestamp{01:18:23--01:19:28}}
\end{classroomqa}

\begin{classroomqa}
\audiencequestion{Is this large-\(N\) approximation what was meant earlier by
coarse-graining?}
\lecturerresponse{No. The thermodynamic limit controls large occupation
numbers and asymptotic counting. Coarse-graining is the loss or deliberate
discarding of microscopic resolution. They are distinct operations, even
though both are often used in macroscopic statistical mechanics.
\lecturetimestamp{01:19:29--01:19:46}}
\end{classroomqa}

\begin{classroomqa}
\audiencequestion{Since
\(\log C=N\log N-\sum_i n_i\log n_i\), should the second term be minimized
rather than maximized?}
\lecturerresponse{Yes; with \(N\) fixed, maximizing \(\log C\) is equivalent
to minimizing \(\sum_i n_i\log n_i\). After substituting \(n_i=NP_i\), the
fixed \(N\log N\) terms cancel and what remains is
\(-N\sum_iP_i\log P_i=NS\). Thus the same operation is described more
naturally as maximizing the positive entropy \(S\).
\lecturetimestamp{01:19:48--01:20:55}}
\end{classroomqa}

\begin{classroomqa}
\audiencequestion{Why not finish the argument immediately by symmetry and set
all the \(n_i\) equal?}
\lecturerresponse{That shortcut ignores the energy labels. The normalization
constraint treats the states symmetrically, but
\(\sum_iP_iE_i=\overline E\) does not: states of different energy contribute
differently. The energy constraint is exactly where the nonuniform
equilibrium distribution enters.
\lecturetimestamp{01:21:15--01:21:41}}
\end{classroomqa}

\section{The Maximum-Entropy Problem}

We can now state the statistical-mechanical rule without ambiguity. Vary the
probabilities \(P_i\) and maximize
\begin{equation}
 S[P]=-\sum_iP_i\log P_i
 \label{eq:sm3-maximize-entropy}
\end{equation}
subject to
\begin{equation}
 \sum_iP_i=1,
 \qquad
 \sum_iP_iE_i=\overline E.
 \label{eq:sm3-maxent-constraints}
\end{equation}
The maximizing probabilities give the sharply preferred occupation fractions
of a large replica ensemble. The occupation numbers fluctuate, but their
relative deviations shrink as \(N\) grows, so the ratios \(n_i/N\) become
well-defined thermodynamic probabilities.

\begin{classroomqa}
\audiencequestion{Are the \(E_i\) fixed state energies while the constraint
uses a separately given total energy?}
\lecturerresponse{Yes. The spectrum \(\{E_i\}\) belongs to the subsystem.
The complete replica collection has a fixed total energy, but energy
continually moves among replicas. Dividing that total by \(N\) fixes the mean
energy \(\overline E\) of one replica, not its energy in each microscopic
realization.
\lecturetimestamp{01:26:12--01:27:07}}
\end{classroomqa}

The problem is easy to state and not yet solved. It is an extremization with
two constraints, so one more mathematical tool is needed.

\section{Lagrange Multipliers}

For an unconstrained function \(F(x_1,x_2,\ldots)\), a stationary point
satisfies
\begin{equation}
 \frac{\partial F}{\partial x_i}=0
 \qquad\text{for every }i.
 \label{eq:sm3-unconstrained-stationary}
\end{equation}
Now require in addition that
\begin{equation}
 G(x_1,x_2,\ldots)=0.
 \label{eq:sm3-one-constraint}
\end{equation}

\begin{classroomqa}
\audiencequestion{What does \(x\) mean if the picture has coordinates \(x\)
and \(y\)?}
\lecturerresponse{Here \(x\) is shorthand for the whole coordinate vector
\((x_1,x_2,\ldots)\). In the two-dimensional drawing its components are
simply \(x\) and \(y\). Later the components will be the probabilities
\(P_1,P_2,\ldots\).
\lecturetimestamp{01:30:12--01:30:20}}
\end{classroomqa}

Geometrically, the level curves \(F=\text{constant}\) form a contour map,
while \(G=0\) is the curve on which the system is allowed to move. A
constrained extremum occurs where one contour of \(F\) is tangent to the
constraint curve. At that point the gradients are parallel:
\begin{equation}
 \nabla F=-\lambda\nabla G
 \label{eq:sm3-parallel-gradients}
\end{equation}
for some number \(\lambda\).

\begin{figure}[t]
\centering
\includegraphics[width=0.88\textwidth]{lecture_03_figure_07.png}
\caption{The board's contour geometry and the augmented function
\(F'=F+\lambda G\). A constrained extremum occurs at tangency.}
\lectureframe{01:37:10}
\end{figure}

Equation \eqref{eq:sm3-parallel-gradients} can be obtained by defining the
augmented function
\begin{equation}
 F'=F+\lambda G
 \label{eq:sm3-augmented-function}
\end{equation}
and demanding ordinary stationarity:
\begin{equation}
 \frac{\partial F'}{\partial x_i}=0.
 \label{eq:sm3-augmented-stationarity}
\end{equation}
The multiplier \(\lambda\) is initially unknown. That does not leave the
system underdetermined, because the original constraint \(G=0\) supplies one
additional equation.

\begin{classroomqa}
\classquestion{Why does adding a term proportional to \(G\) make the
constrained point stationary?}
\lecturerresponse{At the constrained extremum, motion tangent to \(G=0\)
already produces no first-order change in \(F\). The only remaining slope is
normal to the constraint. Since \(\nabla G\) points in that same normal
direction, a suitable \(\lambda\) cancels it:
\(\nabla F+\lambda\nabla G=0\).
\lecturetimestamp{01:35:06--01:35:32}}
\end{classroomqa}

\section{A Two-Variable Example}

Minimize
\begin{equation}
 F(x,y)=\frac{x^2+y^2}{2}
 \label{eq:sm3-example-f}
\end{equation}
subject to
\begin{equation}
 x+y=1.
 \label{eq:sm3-example-constraint}
\end{equation}
Write the constraint as \(G=x+y-1=0\) and form
\begin{equation}
 F'=\frac{x^2+y^2}{2}+\lambda(x+y-1).
 \label{eq:sm3-example-augmented}
\end{equation}
The stationary equations are
\begin{align}
 \frac{\partial F'}{\partial x}&=x+\lambda=0,\\
 \frac{\partial F'}{\partial y}&=y+\lambda=0.
 \label{eq:sm3-example-derivatives}
\end{align}
Thus \(x=y=-\lambda\). The remaining equation is the constraint,
\[
 x+y=-2\lambda=1,
\]
so
\begin{equation}
 \lambda=-\frac12,
 \qquad
 x=y=\frac12.
 \label{eq:sm3-example-solution}
\end{equation}
The point \((1/2,1/2)\) is where the smallest circle centered at the origin
touches the line \(x+y=1\).

\begin{figure}[t]
\centering
\includegraphics[width=0.88\textwidth]{lecture_03_figure_08.png}
\caption{The completed Lagrange-multiplier calculation for minimizing
\((x^2+y^2)/2\) subject to \(x+y=1\).}
\lectureframe{01:44:10}
\end{figure}

With several constraints, introduce one multiplier for each:
\begin{equation}
 F'=F+\lambda_1G_1+\lambda_2G_2+\cdots.
 \label{eq:sm3-many-multipliers}
\end{equation}
Then solve
\begin{equation}
 \frac{\partial F'}{\partial x_i}=0
 \quad\text{for every }i,
 \qquad
 G_a=0
 \quad\text{for every constraint }a.
 \label{eq:sm3-many-equations}
\end{equation}
The lecture proposes a deliberately artificial practice problem: extremize a
quadratic function of \(x,y,z\) while imposing two independent linear
constraints, such as \(x+y=1\) and \(z+x=1\). Its purpose is not the particular
answer but the rule that each constraint brings both one multiplier and one
equation.

\section{The Problem Ready for the Next Lecture}

For statistical mechanics, the coordinates are the probabilities
\(\{P_i\}\), the function to be maximized is the entropy, and the constraints
are normalization and fixed mean energy. Define
\begin{align}
 S[P]&=-\sum_iP_i\log P_i,\\
 G_1[P]&=\sum_iP_i-1,\\
 G_2[P]&=\sum_iP_iE_i-\overline E.
 \label{eq:sm3-final-functions}
\end{align}
The augmented entropy may therefore be written
\begin{equation}
 \Phi[P]
 =S[P]+\alpha G_1[P]+\beta G_2[P].
 \label{eq:sm3-final-augmented}
\end{equation}
The equations to solve are
\begin{equation}
 \frac{\partial\Phi}{\partial P_i}=0
 \quad\text{for every }i,
 \qquad
 G_1=0,
 \qquad
 G_2=0.
 \label{eq:sm3-final-problem}
\end{equation}

\begin{figure}[H]
\centering
\includegraphics[width=0.88\textwidth]{lecture_03_figure_09.png}
\caption{The closing board: entropy as a function of the probabilities,
together with normalization and fixed mean energy.}
\lectureframe{01:53:05}
\end{figure}

At this point the problem is completely posed. Solving
\eqref{eq:sm3-final-problem} produces the equilibrium distribution. The
multiplier conjugate to energy will become the inverse temperature
\(\beta=1/T\), up to the sign convention chosen in
\eqref{eq:sm3-final-augmented}. The next lecture carries out that
differentiation and obtains the Boltzmann weights.
